\begin{table}[h]
\centering
\small
\caption{Draft SAE concept labels produced by the judge-gated labeling pipeline (\texttt{PROMPT\_ORDER=A}, all-opus agents). Each label is the judge's synthesis after round-by-round refinement against held-out contrastive rows. The \emph{macro acc} column is validator accuracy on 50 held-out rows per feature (5 activating + 5 non-activating per dataset). Labels are flagged as drafts pending polarity-inversion fixes on $f_{86}$ and $f_{92}$.}
\label{tab:concept_labels}
\begin{tabular}{l c p{10cm} c}
\toprule
TFM & Concept & Label & macro acc \\
\midrule
Mitra & $f_{6}$ & Activating rows place tail-mass values (rare categorical codes or extreme-percentile numerics) on a localized column subset while contrasts hold baseline-frequency values at those positions. & 0.66 \\
Mitra & $f_{11}$ & Activating rows present a stable shape-level fingerprint on a localized column subset paired with confident predictions; contrasts share the skeleton but break one component. & 0.66 \\
Mitra & $f_{36}$ & Activating rows densely cover a broad column block at mid-band percentiles or rare-frequency categorical levels; contrasts thin out or push isolated positions to extreme tails. & 0.68 \\
Mitra & $f_{86}$ & Activating rows concentrate tail mass on a localized numeric column subset at upper-mid-to-tail percentiles; contrasts lack that subset or invert its polarity. & 0.20 \\
Mitra & $f_{92}$ & Activating rows concentrate tail-mass on a narrow recurring column subset with near-saturated prediction confidence, while contrasts disperse comparable tail-mass across broader non-overlapping columns at lower confidence. & 0.42 \\
\bottomrule
\end{tabular}
\end{table}
